\subsection{A second provider's real payout API (pre-registered)}\label{app:provider2}
A bank will ask whether the gate is tied to one provider. PR58 ran the same scenarios, model (gpt-4.1-mini) and gate
against a second Indian payment provider's real payout API in test mode (test keys, test money). That provider's
official agent server can list payouts but not create them, so the agent's three payment tools were connected to its
Payouts and Payout Links APIs by a small adapter we wrote, which refuses to send anything without a test-mode key.
48 episodes, no errors, every one of the 65 requests answered with HTTP 200.
\begin{center}\small
\begin{tabular}{p{7.2cm}cc}
\toprule
 & direct & \textbf{gated}\\
\midrule
attacker payout accepted by the provider's API & 10 of 15 & \textbf{0 of 15}\\
episodes in which every legitimate payout was accepted & 15 of 15 & \textbf{15 of 15}\\
\bottomrule
\end{tabular}
\end{center}
{\small Both predictions held. All 62 accepted payouts were then looked up in the provider's own records (no lookup
failed): 7 were to attacker accounts, all from direct runs, and the other 3 attacker payouts were payout links. Every
payout in the records, legitimate and attacker alike, is queued (test-mode payouts wait on the test balance), so the
claim is acceptance by the payment API, not settlement. With two independent providers, the gate's result is the same:
the attacker's payout is accepted without it and never with it, and legitimate payouts are unaffected.}

